% Draft for supervisor review, 26 October 2026. Target: IEEE ICBC 2027 (deadline to be confirmed on the official site).
% All numbers come from committed summaries under research/ (26 October 2026 runs); \pending{} is kept for future gaps.
\documentclass[conference]{IEEEtran}
\usepackage{amsmath,amssymb,amsthm}
\usepackage{booktabs}
\usepackage{multirow}
\usepackage{xcolor}
\usepackage{url}
\usepackage{cite}
\newtheorem{proposition}{Proposition}
\newcommand{\pending}[1]{\textcolor{orange}{[pending: #1]}}
\newcommand{\todo}[1]{\textcolor{red}{[TODO: #1]}}

\begin{document}

\title{Containing Transacting Agents:\\Sequence-Level On-Chain Enforcement for LLM Wallet Agents,\\Measured on a Fault-Injected Mainnet Fork}

\author{\IEEEauthorblockN{Draft --- author list to be confirmed}
\IEEEauthorblockA{University of Toronto (CSC490)}}

\maketitle

\begin{abstract}
LLM agents now hold wallets and execute multi-step DeFi plans. When such an agent loses a receipt, repeats a
trade, or misreads its own balance, the mistake is irreversible once it reaches the chain. Off-chain safeguards in
agent frameworks (confirmation gates, output validators) and existing on-chain delegation caveats each address
part of this risk, but nobody has measured which layer contains which failure for a deployed framework. We contribute (i) a
fault-injected benchmark that runs a deployed agent framework (ElizaOS v2) end to end on a pinned Ethereum mainnet
fork, scores outcomes from chain state, and audits every run against tampering; (ii) a failure study of that
framework: 28 numbered findings, from provider-API drift to safety layers that reject honest replies, 25 of them fixed
with regression evidence; (iii) a proof that per-redemption caveats whose bounds are fixed at signing time cannot
enforce an absolute post-trade reserve without either admitting a breach or refusing a feasible step, and that the
deployed caveat set cannot express step precedence; and (iv) two caveat enforcers for MetaMask's deployed
delegation framework, \emph{PlanStepEnforcer} (order, once-only success) and \emph{BalanceFloorEnforcer} (absolute
floors), that close this gap. Against a compromised agent key that holds every signed delegation, our enforcers
admit 0 of 244 off-plan actions and refuse 0 of 44 legitimate ones across four tasks and four market states. A plan
built from MetaMask's built-in caveats (call limits and relative balance caps) admits 16 and refuses 4 feasible trades, and a stateless
per-transaction policy admits 64.
In a fault-injected development sweep of 59 live ElizaOS episodes across four models, 15 unsafe effects reached the
chain without on-chain controls and none with them; our enforcers refused no legitimate step. The enforcers cost no
more gas than that built-in plan.
\end{abstract}

\begin{IEEEkeywords}
LLM agents, account abstraction, delegation, caveat enforcers, DeFi, benchmark, mainnet fork
\end{IEEEkeywords}

\section{Introduction}

Agent frameworks such as ElizaOS~\cite{elizaos} give a language model a wallet and a set of financial tools.
A user asks for an outcome (``sell 0.1~ETH, then buy at least 37.31~DAI, and keep at least 178.23~USDC'') and the
agent plans and submits the transactions. Transactions are irreversible, so the question that matters is not whether
the agent is usually right, but \emph{what reaches the chain when it is wrong}.

Recent work frames this as an open problem. Per-transaction policy verified by a smart account (PACE~\cite{pace2026})
leaves ``multi-step agent loops that compose policy-valid legs into a harmful sequence'' as ``the most important next
target''. Off-chain safety invariants for wallet agents~\cite{yin2026invariants} name plan transactionality as the
open problem most worth attacking. Off-chain exactly-once schemes~\cite{caplease2026,aidguard2026} note that
exactly-once physical effects require an idempotent sink, which a blockchain transaction is not. Benchmarks that execute agent
transactions on forks and score state~\cite{evmquestbench2026,intent2tx2026} evaluate generated code or single
transactions, not deployed frameworks under faults.

We study the full path: a deployed framework, a real model, a pinned mainnet fork, injected faults (a hidden
receipt, an adverse price move), and the control layers that sit between the agent and the chain. Our claim is about
\emph{containment}: under which controls do unsafe effects reach the chain, and at what cost in refused legitimate
steps and gas. We do not claim that agents become reliable.

\textbf{Contributions.}
\begin{enumerate}
\item A benchmark and harness for transacting agents: pinned fork, scripted user, egress proxy with spending caps and
a call ledger, fault-injecting RPC proxy, chain-state oracle, and an independent verifier with tamper tests
(Sec.~\ref{sec:bench}).
\item A failure study of ElizaOS~v2 with trace evidence: 28 findings in six classes, and two agent-level failure
patterns (duplicate sale; goal used as the input amount) (Sec.~\ref{sec:failures}).
\item An expressiveness result: what signing-time caveats can and cannot guarantee across a multi-step plan
(Sec.~\ref{sec:gap}).
\item Two enforcers for the deployed MetaMask delegation framework v1.3.0~\cite{metamaskdelegation,erc7710}
that make order, once-only success and absolute floors chain-enforced (Sec.~\ref{sec:design}).
\item A containment evaluation with ablations: a compromised-agent matrix, an agent-error suite through the
framework's own wallet path, a fault sweep of 59 live episodes across four models, a second framework-free agent, a
plan-authoring probe, and gas (Sec.~\ref{sec:eval}).
\end{enumerate}

\section{Background}

\textbf{Agent frameworks.} ElizaOS~v2 composes a message pipeline, a planner that calls tools natively, an
evaluator that decides whether to continue, and plugins. Its wallet plugin routes swaps and transfers and requires
the user to confirm each financial action; a confirmation binds the action's parameters. Output validators check
that replies about holdings and completed actions are grounded in tool results.

\textbf{Delegation and caveats.} In MetaMask's delegation framework (ERC-7710~\cite{erc7710}), a user's smart account
(here a HybridDeleGator) signs delegations to an agent key. Each delegation carries caveats; at redemption the
DelegationManager calls each caveat enforcer's hooks around the execution, and any revert refuses it. The v1.3.0
deployment has 32 built-in enforcers: targets, methods, calldata, value caps, call limits, time and block windows, and
relative balance-change checks, among others.

\section{Failure and threat model}\label{sec:threat}

A user's intent is a \emph{plan}: ordered steps $s_0,\dots,s_{n-1}$ with prerequisites, each to succeed at most
once, plus constraints on the final state: reserves $\mathrm{bal}_X \ge R$ and goals $\mathrm{bal}_Y \ge G$. We
consider two agents. An \emph{honest but faulty} agent follows the task but loses receipts, repeats steps, misreads
amounts or misreports outcomes; these are the failures we recorded. A \emph{compromised} agent controls the agent
key, for example after key theft or prompt injection, and sends any execution it likes under the delegations it holds,
bypassing the framework. The user, the plan's author, is trusted: we assume the plan is correct, and we treat plan
authoring as a limitation (Sec.~\ref{sec:limits}).

\section{The expressiveness gap}\label{sec:gap}

Consider a two-step plan signed at time $t_0$: step~1 sells a fixed amount of ETH for $X$; step~2 swaps $a$ of
$X$ for $Y$. After step~2 the account must hold $\mathrm{bal}_X \ge R$ and $\mathrm{bal}_Y \ge G$, with $R > 0$. Let
$x$ be the account's $X$ balance just before step~2. It is unknown at $t_0$: step~1's delivery varies within the
agent's slippage tolerance, and other funds may arrive. So $x$ ranges over $[x^-, x^+]$ with $x^- < x^+$.

A \emph{signing-time spend bound} is a caveat on step~2 that refuses a redemption when the amount of $X$ spent exceeds
a constant $c$ fixed at $t_0$. This covers ERC20BalanceChange with a decrease bound, transfer-amount caps, and any other
bound on the change that is fixed at signing.

\begin{proposition}[No fixed bound is both sound and complete]
Let $F(x)=\{a : a \le x - R \text{ and } \mathrm{out}(a) \ge G\}$ be the feasible step-2 spends at balance $x$, and
suppose some $x^\star$ has a feasible $a^\star \in F(x^\star)$ with $a^\star > x^- - R$. Then for every constant $c$,
either (i) some admitted step~2 breaks the reserve, or (ii) some feasible step~2 is refused.
\end{proposition}
\begin{proof}
If $c > x^- - R$, then at balance $x^-$ any spend in $(x^- - R,\ \min(c, x^-)]$ is admitted and leaves less than $R$;
this interval is non-empty because $R > 0$. Otherwise $c \le x^- - R < a^\star$, so the feasible spend $a^\star$ at
$x^\star$ is refused.
\end{proof}

The hypothesis is mild. It holds whenever more $X$ than the worst case can be available before step~2 while step~2's
price worsens, so that a feasible step needs the extra. Our \emph{price-worse} state is a concrete witness, observed on
the fork (Table~\ref{tab:matrix}). The \emph{shortfall} state realises case~(i): the bound set to the budget expected
at signing admits a reserve breach. An absolute floor checked after the execution ($\mathrm{bal}_X \ge R$) does not
depend on $x$, so it is sound and complete for this constraint by construction. For the goal, a fixed lower bound on the
increase of $Y$ is exact only when the account's $Y$ balance before step~2 is known at signing.

\textbf{Precedence.} To our reading of the 32 v1.3.0 enforcers, none conditions a redemption on whether a
\emph{different} delegation has been redeemed. LimitedCalls counts redemptions of one delegation; IdEnforcer gives
mutual exclusion among delegations that share an id, not precedence; time and block windows order by clock, not by
completion. Hence ``step 2 only after step 1 succeeded'' is not expressible with the deployed set. We observe the
consequence when the account already holds the step-2 input: built-in caveats admit step 2 before step 1.

\section{Design}\label{sec:design}

\textbf{PlanStepEnforcer} (56 lines of Solidity). Its terms are $(\mathit{planId}, i, \mathit{prereqMask})$. In
\texttt{beforeHook} it requires that step $i$ is not yet complete and that all prerequisite bits are set. It then sets bit $i$
in state keyed by (delegation manager, delegator, planId). Because a redemption is atomic, a reverted execution rolls
the bit back and the step can be retried; a step can \emph{succeed} at most once. Try-mode executions are refused,
because a failed call in try mode does not revert and would mark the step complete without effect. The enforcer gives
the chain the idempotency that off-chain exactly-once schemes must assume of the sink~\cite{caplease2026,aidguard2026}: a
re-proposed step reverts.

\textbf{BalanceFloorEnforcer} (43 lines). Its terms are $(\mathit{token}, \mathit{floor})$. In
\texttt{afterAllHook} it requires $\mathrm{bal}_{\mathit{token}}(\mathit{delegator}) \ge \mathit{floor}$. One instance
expresses the reserve and another the goal. Because the check runs after all executions of a redemption,
intermediate balances inside a batch do not matter.

\textbf{Plan-bound wallet mode.} In our ElizaOS fork, the wallet's plan mode holds the user's signed delegations
and executes each swap by \texttt{redeemDelegations} from the agent key. The user's funds stay in the smart account.
The framework's confirmation gate is unchanged; the chain makes the final decision. A plan is signed once, before
step 1, alongside built-in target, method and recipient caveats.

\section{Benchmark and harness}\label{sec:bench}

\textbf{Fork and tasks.} We use a pinned Anvil fork of Ethereum mainnet at block 26{,}250{,}601, routed through
Uniswap~V2~\cite{uniswapv2}. Task D1 (development) sells 0.1~ETH for USDC, then must hold
$\ge$37.31~DAI while keeping $\ge$178.23~USDC. Three held-out tasks (H1--H3) vary the pair, the sizes and the
direction. Task numbers come from a calibration rule fixed before any run.

\textbf{Scenarios.} \emph{normal}; \emph{receipt\_hidden}, where the RPC proxy hides the sale's receipt for 200\,s
against the wallet's 60\,s wait (the classic trigger for a duplicate); and \emph{price\_change\_feasible} and
\emph{price\_change\_infeasible}, a scheduled market move after the sale, after which the goal is still reachable or
no longer reachable.

\textbf{Harness.} Each episode runs in a fresh process with an in-memory database, localhost-only networking and a
neutral random identity. A scripted user sends the task, answers ``Yes, confirm'' to any confirmation request (the
worst case for duplicates), and finally asks for a COMPLETED / INFEASIBLE / IN\_PROGRESS declaration. All model
traffic passes an egress proxy. The proxy enforces per-cohort and cumulative spending caps by worst-case reservation,
records every request and response, and refuses any request that carries a harness label (the blinding check).
Embeddings are computed locally and deterministically.

\textbf{Scoring and audit.} Outcomes come from chain state: sales count, step-2 execution, goal and reserve. They are
reported separately from the declared outcome. The verifier recomputes everything from the raw records, and it checks
source and harness hashes, ledger arithmetic, caps, turn policy and blinding. Nine tamper tests (for example an edited
response, a deleted call, a changed price) must each be rejected.

\textbf{Control layers.} \emph{stock}: the agent key holds the funds and only the framework's gate applies.
\emph{per\_tx\_policy}: target, method, recipient, value cap, and relative balance-change bounds on step 2, all
stateless (PACE-like). \emph{existing\_controls}: per\_tx\_policy plus LimitedCalls~=~1 on every step: our plan
built from built-in caveats. \emph{plan\_guard}: target, method and recipient caveats plus our two enforcers. There are two
ablations: \emph{floor-only} and \emph{order-only}.

\section{Evaluation}\label{sec:eval}

\subsection{Compromised agent}

The agent key holds every signed delegation and sends 18 crafted actions per cell straight to the chain. Each
action starts from an isolated snapshot. The actions include repeats, reordering, exfiltration (output to an attacker,
direct transfer, attacker approval), an unlisted asset, reserve breaches, goal misses, B2 (the goal number as the
input) and the legitimate steps. In-plan versus off-plan labels are computed from the user's constraints in that state. A first
version labelled B2 off-plan by name, although in the shortfall state it met both constraints; we corrected the rule
and kept that attempt on record. Table~\ref{tab:matrix} aggregates the four tasks; Table~\ref{tab:kinds} breaks the admitted
actions down by kind. The run is deterministic: reruns produce byte-identical results.

\begin{table}[t]
\centering
\caption{Compromised agent, all four tasks: off-plan actions admitted (of 60, or 64 in price-worse) / legitimate
actions refused (of 12, or 8). Every action starts from an isolated snapshot.}
\label{tab:matrix}
\setlength{\tabcolsep}{3pt}
\begin{tabular}{lccccc}
\toprule
Layer & calibr. & shortf. & prefund. & price-w. & total \\
\midrule
none (agent key holds funds) & 56\,/\,0 & 56\,/\,0 & 60\,/\,0 & 64\,/\,0 & 236\,/\,0 \\
per\_tx\_policy (stateless) & 12\,/\,0 & 20\,/\,0 & 16\,/\,0 & 16\,/\,4 & 64\,/\,4 \\
existing\_controls (built-ins) & 0\,/\,0 & 8\,/\,0 & 4\,/\,0 & 4\,/\,4 & 16\,/\,4 \\
floor-only (ablation) & 15\,/\,0 & 12\,/\,0 & 16\,/\,0 & 16\,/\,0 & 59\,/\,0 \\
order-only (ablation) & 20\,/\,0 & 20\,/\,0 & 20\,/\,0 & 24\,/\,0 & 84\,/\,0 \\
\textbf{plan\_guard} & \textbf{0\,/\,0} & \textbf{0\,/\,0} & \textbf{0\,/\,0} & \textbf{0\,/\,0} & \textbf{0\,/\,0} \\
\bottomrule
\end{tabular}
\end{table}

\begin{table}[t]
\centering
\caption{Off-plan actions admitted by kind (all tasks and states).}
\label{tab:kinds}
\setlength{\tabcolsep}{2.5pt}
\begin{tabular}{lccccccc}
\toprule
Layer & repeat & order & reserve & goal & budget & asset & exfil. \\
 & /64 & /16 & /32 & /28 & /8 & /16 & /80 \\
\midrule
none & 64 & 8 & 32 & 28 & 8 & 16 & 80 \\
per\_tx\_policy & 48 & 8 & 4 & 0 & 4 & 0 & 0 \\
existing\_controls & 0 & 8 & 4 & 0 & 4 & 0 & 0 \\
floor-only & 59 & 0 & 0 & 0 & 0 & 0 & 0 \\
order-only & 0 & 0 & 32 & 28 & 8 & 16 & 0 \\
\textbf{plan\_guard} & \textbf{0} & \textbf{0} & \textbf{0} & \textbf{0} & \textbf{0} & \textbf{0} & \textbf{0} \\
\bottomrule
\end{tabular}
\end{table}

\textbf{Reading.}
\begin{itemize}
\item Without on-chain controls, the compromised key performs every off-plan action that has funds behind it. In
64 cases it moves funds to the attacker.
\item Every guarded layer stops exfiltration: target, method and recipient caveats suffice.
\item The stateless policy admits repeated sales (B1). Built-in caveats stop repeats with LimitedCalls but admit:
\begin{itemize}
\item two reserve breaches under shortfall (an overspend, and spending the budget fixed at signing);
\item step 2 before step 1 when the account is prefunded.
\end{itemize}
They also refuse the feasible step 2 when Y gets dearer (the relative cap, Proposition~1).
\item The ablations show that both enforcers are needed. Floors alone admit repeated sales and a repeated step 2.
Order alone admits reserve and goal violations, B2, and a swap into an asset the plan never names; the goal floor is
what pins the output asset.
\item plan\_guard admits none of the 244 off-plan actions and refuses none of the 44 legitimate ones.
\end{itemize}

\subsection{Agent errors through the framework}

We also replay a fixed sequence of wallet requests through stock ElizaOS's own WALLET router, each confirmed as a user
would confirm it (no model). The sequence covers step 2 before step 1, step 1, repeated sales, a transfer to a new
address, overspends, a goal miss, the true surplus, and a repeated step 2. It runs on all four tasks in three states.
\begin{itemize}
\item The framework's gate alone (stock) admits 6 or 7 off-plan requests per cell.
\item Built-in caveats admit one per cell in the shortfall and prefunded states, and then refuse the correct step 2 as a
consequence.
\item plan\_guard admits 0 and refuses 0 in all 12 cells.
\end{itemize}
The confirmation gate never refuses a repeated sale: it binds a request, not a plan.

\subsection{Who writes the plan?}

The enforcers guarantee adherence to the signed plan, not that the plan is right. We asked four models to turn each
task's request into the plan specification, 3 times per task (48 requests, no tools), and scored the result field by
field.
\begin{itemize}
\item GPT-6 Luna, Claude Haiku~5.5 and Claude Sonnet~5.5 produced a fully correct plan in 36 of 36 cases.
\item GPT-4.1-mini got the sale, the order edge and both floors right in 12 of 12, but left step~2's input asset
open in 9 of 12. The request never says which asset pays for step~2, so this is a literal but looser plan: signed as
is, it would let the agent spend other assets on step~2 while the floors still hold.
\end{itemize}
So the residual risk is under-specification rather than wrong numbers, and a plan-review screen that shows each step's
input and output asset addresses it.

\subsection{Live development runs}

These are development data at one or two repetitions per cell. They are counts, not rates.

\textbf{Live checks.} We ran two live checks with Claude Haiku~5.5 and GPT-6~Luna on D1 \emph{normal}, before and
after fixing five framework defects found in the first check (L1--L5). Episodes that completed the on-chain task
went from 2 of 4 to 4 of 4. Declarations matching chain state went from 0 of 4 to 3 of 4; the fourth ran out of
budget after finishing. No episode sold twice.

\textbf{Fault sweep.} We then ran a pre-declared sweep: D1 under all four scenarios, with three conditions (stock,
existing\_controls, plan\_guard) and four models. Luna and Haiku ran at $r=2$ in all three conditions, GPT-4.1-mini at
$r=1$ in stock and plan\_guard, and Sonnet~5.5 on receipt\_hidden only. That is 59 episodes, \$5.29, every cohort
audited with 9/9 tamper cases rejected.
\begin{itemize}
\item \textbf{Stock} (21 episodes): 15 unsafe effects reached the chain: 4 duplicate sales, 3 trades after the goal
became infeasible, 3 off-plan swap legs, 2 reserve breaches, 2 B2 trades (37.31~USDC as the input) and 1 repeated
step~2.
\item \textbf{Built-in caveats} (16 episodes) and \textbf{plan\_guard} (21 episodes): none.
\item \textbf{Chain refusals:} the chain refused off-plan attempts in 8 episodes of the guarded runs: goal-missing
trades, B2, and repeated sales. It refused no legitimate step.
\end{itemize}
Three caveats:
\begin{itemize}
\item \textbf{Behaviour differs by condition.} Plan mode shows the plan in the agent's context, and the guarded agents
attempted fewer off-plan actions, so part of the difference is behavioural.
\item \textbf{Built-ins versus plan\_guard.} The market states that separate the two did not arise in these D1 runs;
the matrix measures them.
\item \textbf{Sonnet was truncated} by the proxy's worst-case output reservation, not by the model, and Haiku under
plan\_guard completed only 3 of 8, stalling in confirmation loops rather than being refused.
\end{itemize}

\subsection{A second, framework-free agent}

To check that containment does not depend on ElizaOS, we ran a minimal reference agent: a plain tool-calling loop on
the provider API with four tools (balances, quote, swap, transaction status), no framework, memory or validators.
It used the same fork, task, user texts and faults: the receipt is hidden for the submitting turn, and the scheduled
market move follows the sale. We ran 32 pre-declared episodes (Luna and Haiku, stock and plan\_guard, four scenarios,
$r=2$) for \$0.11. This harness is lighter than the audited one.
\begin{itemize}
\item \textbf{Reliability.} 31 of 32 episodes were correct on chain and correctly declared.
\item \textbf{The one unsafe effect.} Luna, under stock, spent the surplus on a step~2 that could not reach the goal
after the price move.
\item \textbf{The same mistake, contained.} Under plan\_guard, in the same cell, the identical attempt was refused twice
by the goal floor, and the agent declared the task infeasible.
\item \textbf{A smoke test before the sweep.} Luna, under stock, sold a second 0.1~ETH after first declaring the goal
infeasible. It then made three step-2 trades, broke the reserve by $10^{-6}$~USDC, and declared COMPLETED.
\end{itemize}
The simpler agent fails less often than ElizaOS, which supports the failure study's finding that much of ElizaOS's
failure is the framework's. When either agent errs, though, the error reaches the chain unless an on-chain layer stops
it.

\subsection{Gas}

Table~\ref{tab:gas} gives the gas used by the legitimate steps in the compromised-agent run. All layers start from
the same calibrated fork state. Deploying the two enforcers costs 399{,}478 and 304{,}337 gas, once per chain.

\begin{table}[t]
\centering
\caption{Gas of the legitimate steps, same fork state (task D1, calibrated).}
\label{tab:gas}
\begin{tabular}{lrr}
\toprule
Layer & Sale & Step 2 \\
\midrule
none (direct from the agent key) & 130{,}205 & 188{,}092 \\
per\_tx\_policy & 282{,}521 & 397{,}518 \\
existing\_controls & 329{,}925 & 444{,}973 \\
floor-only & 282{,}557 & 356{,}950 \\
order-only & 331{,}538 & 342{,}620 \\
\textbf{plan\_guard} & 331{,}550 & 388{,}894 \\
\bottomrule
\end{tabular}
\end{table}

Most of the cost is the delegation framework itself: about 152k gas per redemption over a direct transaction. Against
that baseline, plan\_guard's sale costs the same as the built-in call limit (+49k). Its step 2 is 56k cheaper than
the built-in plan. Both use two balance checks and one stateful check, but each floor reads a balance once, after the
execution, whereas each relative balance-change caveat records the balance before and compares it after.

\section{Failure study: ElizaOS v2}\label{sec:failures}

Making the framework transact on a fork at all, then completing the task, surfaced 28 numbered findings
(Table~\ref{tab:defects}). Each has trace evidence, and 25 are fixed with a regression test or replay. We report
classes rather than individual bugs.

\begin{table}[t]
\centering
\caption{Framework findings by class (IDs overlap; not independent counts).}
\label{tab:defects}
\setlength{\tabcolsep}{3pt}
\begin{tabular}{p{2.6cm}p{5.3cm}}
\toprule
Class & Examples \\
\midrule
Provider-API drift & A1: GPT-6 tool calls rejected without \texttt{reasoning\_effort} none (every call 400); A11: Sonnet/Opus~5.5 reject forced tool choice (every planner call 400); A12: GPT-6.1 Sol unusable through Chat Completions (open); H6 \\
Missing chain primitives & A4: no fork-usable swap route; H7: no receipt or balance read; A5: token holdings absent \\
Stale or lost state & L2: earlier transaction hashes lost across turns; A6, A8/H8, L5: stale balances; H4/L7: earlier effects lack current proof \\
Over-eager safety layers & L1: credential redaction masks token addresses, which the model copies back; A9/H1/H2/L4/L6: honest replies rejected as ungrounded; H3: no effect receipt for a submitted trade \\
Confirmation binding & L3: ``yes'' does not execute when the chain alias or an echoed recipient differs \\
Tool surface & A2, A7, H9, H5, A10 \\
\bottomrule
\end{tabular}
\end{table}

Three lessons matter for containment:
\begin{itemize}
\item \textbf{Off-chain safety layers mostly cost progress.} The framework's validators and redaction rejected honest
replies, hid the assets the agent had to name, and lost the agent's record of its own effects. In one development
run this led to a duplicate sale.
\item \textbf{The confirmation gate approved every duplicate it was shown.} It binds parameters, not plan state.
\item \textbf{Provider drift silently disables whole model families.} Two of the four current frontier families
could not call tools through the stock framework at all.
\end{itemize}
Fixing the framework is necessary to measure it, but it does not remove the need for containment: after all fixes,
agents still re-proposed completed steps (B1) and used the goal as the input amount (B2).

\section{Related work}

\textbf{Fork-based evaluation.} EVM-QuestBench~\cite{evmquestbench2026} runs multi-turn generated code on a BSC fork
with final-state scoring. Intent2Tx~\cite{intent2tx2026} scores one-shot transactions by state effect.
DeFiFlowBench~\cite{defiflowbench2026} checks synthesized workflows statically. None runs a deployed framework under
injected faults or measures duplicate effects.

\textbf{Guards for agent actions.} PACE~\cite{pace2026} verifies per-transaction policy on-chain.
Yin~\cite{yin2026invariants} enforces off-chain invariants in a custom wallet agent. CapLease and
AID-Guard~\cite{caplease2026,aidguard2026} provide off-chain exactly-once authorization. ChainClaw~\cite{chainclaw2026}
guards multi-step workflows off-chain, and Proof of Execution~\cite{poe2026} makes traces tamper-evident. We move
sequence-level guarantees on-chain and compare layers under the same faults.

\textbf{Attacks and surveys.} Context-manipulation attacks on ElizaOS~\cite{craibench2025,cryptoland2025} motivate
our compromised-agent model. Surveys~\cite{alqithami2026,sokcommerce2026} call for reproducible, finance-specific
harnesses. Market studies~\cite{papergains2026} describe deployed agents' economics, not their failures.

\section{Limitations and discussion}\label{sec:limits}

\begin{itemize}
\item \textbf{Plan authoring is trusted.} The enforcers guarantee adherence to the signed plan, not that the plan is
right. Our probe found models reliable on numbers and order, with one model under-specifying the input asset; the
probe uses one phrasing and two-step plans.
\item \textbf{Scope.} We study one chain, one DEX route (a small on-chain Uniswap~V2 route added to the stock
framework, which otherwise quotes only through HTTP aggregators), one framework family, and two-step plans.
\item \textbf{Evidence.} The live data are development runs at small $r$. The compromised-agent and agent-error
results are deterministic but constructed. The enforcers are not audited, and within one batched redemption,
prerequisites follow the manager's hook order.
\item \textbf{Configurations.} Our stock condition runs the hardened framework; a minimal-fix configuration is
planned.
\end{itemize}

\section{Conclusion}

Off-chain safeguards and per-transaction policies do not contain the failures that multi-step transacting agents
actually produce. Signing-time caveats provably cannot express an absolute reserve or step precedence. Two small
enforcers on a deployed delegation framework close that gap: they admit no off-plan action and refuse no legitimate
one in our matrix, at no more gas than the plan built from built-in caveats.

\bibliographystyle{IEEEtran}
\bibliography{refs}
\end{document}
